% !TEX program = lualatex
% Hoja de ejercicios: métodos de demostración y condiciones lógicas

\documentclass[11pt,a4paper]{article}

\usepackage{fontspec}
\usepackage[spanish,es-nodecimaldot]{babel}
\usepackage[a4paper,margin=2.35cm,headheight=15pt]{geometry}
\usepackage{amsmath,amssymb,mathtools}
\usepackage{enumitem}
\usepackage{tabularx}
\usepackage{xcolor}
\usepackage{fancyhdr}
\usepackage{lastpage}
\usepackage{hyperref}

\setmainfont{Latin Modern Roman}
\setsansfont{Latin Modern Sans}
\setmonofont{Latin Modern Mono}

\definecolor{FdMBlue}{RGB}{20,55,100}
\definecolor{FdMLightBlue}{RGB}{235,241,248}
\definecolor{FdMGray}{RGB}{90,96,104}

\hypersetup{
  colorlinks=true,
  linkcolor=FdMBlue,
  urlcolor=FdMBlue,
  pdftitle={Ejercicios: métodos de demostración},
  pdfauthor={Fundamentos de las Matemáticas I}
}

\pagestyle{fancy}
\fancyhf{}
\lhead{Fundamentos de las Matemáticas I}
\rhead{Métodos de demostración}
\cfoot{\thepage\ de \pageref{LastPage}}
\renewcommand{\headrulewidth}{0.4pt}

\setlist[enumerate]{leftmargin=2.4em,itemsep=0.35em,topsep=0.4em}
\setlist[itemize]{leftmargin=2em,itemsep=0.2em,topsep=0.3em}

\newcounter{ejercicio}
\newenvironment{ejercicio}[1][]{%
  \refstepcounter{ejercicio}%
  \par\medskip\noindent
  {\color{FdMBlue}\sffamily\bfseries Ejercicio \theejercicio%
  \if\relax\detokenize{#1}\relax\else\ --- #1\fi.}\quad
}{\par\medskip}

\newcommand{\respuesta}[1]{%
  \par\vspace{#1}\noindent
  \textcolor{FdMGray!45}{\rule{\linewidth}{0.3pt}}\par
}

\newcommand{\cajarecordatorio}[1]{%
  \begin{center}
    \colorbox{FdMLightBlue}{%
      \begin{minipage}{0.91\linewidth}
        \small #1
      \end{minipage}}
  \end{center}
}

\begin{document}

\begin{center}
  {\color{FdMBlue}\sffamily\bfseries\LARGE Hoja de ejercicios}\par
  \vspace{0.2em}
  {\sffamily\large Métodos de demostración y condiciones lógicas}\par
  \vspace{0.7em}
  \rule{\linewidth}{0.6pt}
\end{center}

\noindent
\begin{tabularx}{\linewidth}{@{}lX@{\hspace{1.5em}}lX@{}}
  \textbf{Nombre y apellidos:} & \hrulefill &
  \textbf{Fecha:} & \hrulefill
\end{tabularx}

\vspace{0.8em}
\noindent
En cada prueba deben declararse las variables, las hipótesis y el método
empleado. Cada paso debe justificarse mediante una definición, una propiedad
conocida o una inferencia lógica válida.

\cajarecordatorio{
  \textbf{Definiciones previas.}
  Sean \(a,b\in\mathbb Z\). Decimos que \(b\) es \emph{múltiplo de \(a\)}
  si existe \(k\in\mathbb Z\) tal que \(b=ak\).
  Esta misma relación se escribe \(a\mid b\) y se lee
  \emph{\(a\) divide a \(b\)}. En consecuencia,
  \[
    a\mid b
    \quad\Longleftrightarrow\quad
    b\text{ es múltiplo de }a.
  \]
}

\cajarecordatorio{
  \textbf{Recordatorio lógico.}
  \(P\) es suficiente para \(Q\) cuando \(P\Rightarrow Q\);
  \(P\) es necesaria para \(Q\) cuando \(Q\Rightarrow P\).
  Para probar \(P\Leftrightarrow Q\) hay que demostrar las dos implicaciones.
}

\section*{A. Condiciones necesarias y suficientes}

\begin{ejercicio}[Lectura de una implicación]
Para cada pareja de propiedades, indica cuál es condición suficiente y cuál es
condición necesaria. Decide también si se trata de una condición necesaria y
suficiente.
\begin{enumerate}[label=\alph*)]
  \item Para \(n\in\mathbb Z\): ser múltiplo de \(12\) y ser múltiplo de \(4\).
  \item Para un cuadrilátero: ser un cuadrado y ser un rectángulo.
  \item Para \(n\in\mathbb Z\): ser par y ser de la forma \(2k\) para algún
        \(k\in\mathbb Z\).
\end{enumerate}
\respuesta{4.3cm}
\end{ejercicio}

\begin{ejercicio}[Traducir condiciones a implicaciones]
Sean \(P\) y \(Q\) dos proposiciones. Escribe con símbolos:
\begin{enumerate}[label=\alph*)]
  \item \(P\) es condición suficiente para \(Q\).
  \item \(P\) es condición necesaria para \(Q\).
  \item \(P\) es condición necesaria y suficiente para \(Q\).
\end{enumerate}
Después, si \(P(n)\) significa que \(n\) es múltiplo de \(6\) y \(Q(n)\)
significa que \(n\) es múltiplo de \(3\), describe la relación entre ambas.
\respuesta{4.2cm}
\end{ejercicio}

\newpage

\begin{ejercicio}[Necesaria pero no suficiente]
Para \(x\in\mathbb R\), considera la propiedad
\[
  R(x):\quad x^2>x.
\]
Demuestra que \(x\neq0\) es una condición necesaria, pero no suficiente, para
que se cumpla \(R(x)\). Comprueba también que \(x>1\) es suficiente, pero no
necesaria.
\respuesta{6.0cm}
\end{ejercicio}

\begin{ejercicio}[Una condición necesaria y suficiente]
Demuestra que un entero \(n\) es múltiplo de \(10\) si y solo si es par y
múltiplo de \(5\).
\respuesta{7.0cm}
\end{ejercicio}

\newpage
\section*{B. Demostraciones directas y por contraposición}

\begin{ejercicio}[Demostración directa]
Demuestra directamente que el cuadrado de todo número entero impar es impar.
\respuesta{5.0cm}
\end{ejercicio}

\begin{ejercicio}[Racionales cerrados para la suma]
Demuestra directamente que la suma de dos números racionales es racional.
\respuesta{5.0cm}
\end{ejercicio}

\begin{ejercicio}[Transitividad de la divisibilidad]
Sean \(a,b,c\in\mathbb Z\). Demuestra directamente que, si \(a\mid b\) y
\(b\mid c\), entonces \(a\mid c\).
\respuesta{4.5cm}
\end{ejercicio}

\newpage

\begin{ejercicio}[Contraposición y paridad]
Demuestra por contraposición que, si \(n^2\) es par, entonces \(n\) es par.
\respuesta{7.0cm}
\end{ejercicio}

\begin{ejercicio}[Contraposición y divisibilidad]
Demuestra por contraposición que, si \(3\mid n^2\), entonces \(3\mid n\).
\respuesta{7.5cm}
\end{ejercicio}

\newpage
\section*{C. Contradicción y división en casos}

\begin{ejercicio}[Irracionalidad de \(\sqrt2\)]
Demuestra por contradicción que \(\sqrt2\notin\mathbb Q\).
\respuesta{7.5cm}
\end{ejercicio}

\begin{ejercicio}[No existe el mayor entero]
Demuestra por contradicción que no existe un número entero mayor o igual que
todos los números enteros.
\respuesta{5.5cm}
\end{ejercicio}

\newpage

\begin{ejercicio}[Dos enteros consecutivos]
Demuestra por casos que, para todo \(n\in\mathbb Z\), el número
\[
  n(n+1)
\]
es par.
\respuesta{7.0cm}
\end{ejercicio}

\begin{ejercicio}[Valor absoluto]
Demuestra por casos que, para todo \(x\in\mathbb R\),
\[
  |x|=x\quad\Longleftrightarrow\quad x\geq0.
\]
\respuesta{7.0cm}
\end{ejercicio}

\newpage
\section*{D. Equivalencias, existencia y unicidad}

\begin{ejercicio}[Equivalencia de paridad]
Demuestra que, para todo \(n\in\mathbb Z\),
\[
  n\text{ es par}\quad\Longleftrightarrow\quad n^2\text{ es par}.
\]
Indica el método utilizado en cada implicación.
\respuesta{7.0cm}
\end{ejercicio}

\begin{ejercicio}[Existencia y unicidad]
Sean \(a,b\in\mathbb R\), con \(a\neq0\). Demuestra que existe un único
\(x\in\mathbb R\) tal que
\[
  ax+b=0.
\]
\respuesta{7.0cm}
\end{ejercicio}

\newpage

\begin{ejercicio}[Diagnóstico de una demostración]
Se quiere demostrar:
\[
  n^2\text{ es múltiplo de }4\Rightarrow n\text{ es par}.
\]
Un estudiante escribe: \emph{Si \(n\) es par, entonces \(n=2k\), y por tanto
\(n^2=4k^2\) es múltiplo de \(4\).} Explica por qué esto no demuestra el
enunciado y presenta una demostración correcta.
\respuesta{7.0cm}
\end{ejercicio}

\begin{ejercicio}[Refutación mediante contraejemplo]
Refuta el enunciado
\[
  \forall x\in\mathbb R,\qquad x^2\geq x.
\]
Explica por qué el procedimiento empleado es suficiente.
\respuesta{5.5cm}
\end{ejercicio}

\newpage

\begin{ejercicio}[Elección del método]
Indica un método especialmente adecuado para cada objetivo y justifica
brevemente la elección.
\begin{enumerate}[label=\alph*)]
  \item Si \(n\) es impar, entonces \(n^2\) es impar.
  \item Si \(n^2\) es par, entonces \(n\) es par.
  \item No existe un mayor número entero.
  \item Para todo \(n\in\mathbb Z\), \(n(n+1)\) es par.
  \item \(n\) es par si y solo si \(n^2\) es par.
\end{enumerate}
\respuesta{9.0cm}
\end{ejercicio}

\end{document}
